\chapter{Historical notes}
\label{chap:history}

These notes record the sources of the constructions and results developed in the mathematical chapters. References to numbered results identify the statement used in this book; the accompanying source locators specify the versions and hypotheses against which it was checked. A related result with stronger hypotheses is distinguished from the argument given here.

\section{Plurisubharmonic functions and regularization}
\label{sec:history-psh}

The normalized $L^1$ estimate in \cref{thm:normalized-compactness} is also recorded in \cite[Proposition~6.2]{Dinew16}. The proof here uses local psh compactness and propagation over a connected coordinate cover. The local exponential-integrability input for \cref{lem:uniform-exponential} is recalled in \cite[Section~1]{Zeriahi01}; the weighted-sum contradiction then gives constants uniform over normalized potentials.

The comparison of convolutions in different coordinates in \cref{lem:bounded-local-regularization,cor:coordinate-convolution} is the bounded part of the argument in \cite[\S2, Lemmas~4--5]{BK07}. The full regularization theorem of \cite[Theorem~1]{BK07} permits unbounded potentials and a positive continuous background. The proof of \cref{thm:bounded-exact-regularization} retains the bounded case needed for the product constructions, including the decreasing diagonal and the rescaling that preserves a Hermitian background.

The Chern-connection expansion in \cite[Theorem 4.1]{Dem94} applies to the closed current $\ddc p$ with arbitrary smooth lower bound $-\theta$. Its use for \cref{thm:big-continuous-envelope-contact} retains unbounded potentials with finite $L^1$ norm; radius localization supplies the minimum-principle argument off the singular barrier.

The resolvable regularization of \cref{thm:resolvable-regularization-input} is \cite[Theorem 3.2]{DP04}. It allows an arbitrary smooth lower bound on a Hermitian manifold and supplies a smooth remainder after resolution. This is a separate regularizing sequence from the prescribed multiplier-ideal sequence in \cref{thm:demailly-hermitian-input}; simultaneous possession of their additional properties is not assumed.

The monotonicity and logarithmic convexity used for \eqref{eq:smooth-background-transform-rate} come from \cite[Theorem~4.1]{Dem94}. The proof separately establishes a linear error rate for a smooth lower bound.

The curvature estimate in the proof of \eqref{eq:smooth-background-transform-rate} is the algebraic inequality used in \cite[Lemma~4.4]{Dem94}, here with no small $\epsilon$ term or $M_\epsilon$ correction.

The error estimates for \eqref{eq:smooth-background-transform-rate} use the expansion in \cite[Proposition~3.8, equations (3.11), (4.3), and (4.5)]{Dem94}. Local trace bounds and scaling give constants uniform in the center and averaging radius.

The Chern-connection estimate in \cref{lem:chern-kiselman-input} comes from \cite[Theorem 4.1 and Remark 4.7]{Dem94}; the stated truncated transform is recorded in \cite[Lemma 4.1]{KN19}. Constants absorb the fixed normalization of $\ddc$.

The full base--fiber expansion used in \cref{lem:general-kiselman-modulus} is \cite[Theorem 4.1]{Dem94}, applied to a closed Hessian with a possibly nonclosed smooth lower bound. Its error is retained as a modulus tending to zero.

The linear-radius term in \cref{lem:general-kiselman-modulus} absorbs the mixed and vertical errors as in \cite[Remark 4.7]{Dem94}.

The fixed non-pluripolar measure convergence in \cref{lem:fixed-npp-integrals} uses Cegrell's classical fixed-measure argument, also written in extraction form in \cite[Lemma 2.1 and Corollary 2.2]{KN22}. Applying the extraction statement to each subsequence gives convergence of the whole bounded sequence.

The analyticity of positive Lelong level sets recalled in \cref{subsec:psh-lelong-levels} and used in \cref{lem:lower-mass-trace-concentration} is Siu's theorem \cite[Main Theorem, p.~53]{Siu74}. Smooth local corrections allow its application to quasi-plurisubharmonic functions.

\section{Mass and energy}
\label{sec:history-mass-energy}

Torsion-corrected energy functionals were constructed by Yi Li; see \cite[Equation~(1.18) and Theorem~1.2]{YiLi10}. A correction changes the one-form being integrated. These constructions are therefore consistent with the closedness obstruction for the raw Monge--Amp\`ere one-form in \cref{thm:raw-primitive}. The radial functional in \cref{prop:radial-torsion-defect} is compared with the raw operator by computing its first variation.

The failure of the bounded mass property on a three-dimensional Hopf manifold is due to Xia--Zhang \cite{XZ26}; the construction is developed in \cref{chap:hopf}. In conjunction with \cref{thm:infinite-mass-witness}, it gives a potential with infinite extended mass and trivial multiplier ideals at every positive weight. The coefficient-integrability and convergence statements in \cref{sec:integrable-infinite-mass} describe the linear current carried by that construction.

The Hopf threefold construction of \cref{chap:hopf} is Xia--Zhang's counterexample \cite[Theorem 1.1]{XZ26}. It uses a Gauduchon background and converts vertical Dirichlet energy into Monge--Amp\`ere mass.

The operator in \cite[Section 1]{XZ26} is $\mathrm i\partial\bar\partial$. In \cref{chap:hopf}, all background forms are divided by $2\pi$ and potentials are unchanged. This preserves the scalar positivity inequalities and changes the coefficient of the mass identity.

The two failures emphasized in the preface are the unbounded-mass example of \cite{XZ26} and the vanishing-lower-mass example of \cite{LowerMass26}. Their constructions are developed in \cref{chap:hopf,sec:vanishing-lower-mass}.

The mass-invariance characterization in \cref{thm:mass-invariance-characterization} uses Chiose's polarization argument \cite[Theorems~1--2]{Chiose16}.

The sufficient lower-bound argument in \eqref{eq:threefold-mass-dirichlet} is the dimension-three case of \cite[Theorem~3.2]{AGL23}. The condition $\ddc\omega\leq0$ is called plurinegativity there.

The principal elliptic bundles in \cref{sec:primitive-elliptic-bundles} are a special case of the torus bundles of Goldstein--Prokushkin \cite[Section 2]{GP04}. Their metric calculation gives an upper bound and does not decide whether lower mass vanishes.

The construction in \cref{sec:vanishing-lower-mass} reconstructs the supplied manuscript \cite{LowerMass26}, with the normalization used throughout this book.

\section{Gauduchon metrics and positivity}
\label{sec:history-gauduchon}

Gauduchon's conformal theorem is recalled immediately before \cite[Lemma~2.3]{TW10}, in the notation $\gamma=\mathrm{e}^f\eta$. The multiplier in \cref{thm:gauduchon} is related by $\rho=\mathrm{e}^{(n-1)f}$. The trace comparison in \cref{cor:trace-bounds} is the mechanism behind the first mixed-volume bounds in \cite[Proposition~3.3]{GL22}.

The positivity criterion for an arbitrary smooth background in \cref{thm:lamari} is recorded in \cite[Lemma~1.15]{BGL25}. The proof here combines separation with a scalar elliptic right inverse. Trace estimates of the form \eqref{eq:gauduchon-c-bound} also enter the weak existence arguments in \cite[Section~4]{BGL25}.

\section{Bounded products and extended measures}
\label{sec:history-products}

The local expansion defining bounded products in \cref{prop:bounded-products} is discussed in \cite[Section~1]{KN15} and \cite[Section~2.1]{BGL25}. The restricted-truncation construction in \cref{thm:arbitrary-extended-ma} uses the local mechanism of \cite[Proposition~2.1]{BGL25}. The finite mixed-current construction of \cite[Theorem~2.2]{BGL25} uses a uniform variation bound before removing the restrictions. Allowing an extended positive measure retains the construction when its total mass is infinite, but such a measure is no longer a global current of order zero.

The metric independence and reduction to smooth potentials in \cref{prop:bmp-metric} recover the comparison recorded after \cite[Definition~1.16]{BGL25}. Exact decreasing regularization supplies the passage from smooth to bounded potentials.

The Gauduchon trace bound and cancellation of the exact top wedge used in \cref{prop:gauduchon-range-growth} also appear in \cite[Proposition 3.3(i), (vi)]{GL22}. The range estimate retains their polynomial dependence on the oscillation.

The local mixed inequality in \cref{thm:local-mixed-input} is \cite[Theorem 1.2]{Dinew09}, for locally bounded psh functions and nonnegative measurable smooth-volume densities. Its extension to smooth backgrounds, proved in \cref{sec:psh-bounded-products}, is used in \cref{sec:interior-calculus}. Freezing the two smooth backgrounds gives \cref{lem:nonclosed-mixed-density} and the trace inequalities \eqref{eq:big-singular-gauduchon-trace} and \cref{lem:unbounded-gauduchon-trace-density}; the latter retains the linear trace mass on the pole set. The coefficient-freezing proof also recovers the smooth-volume inequality used in \cite[Lemma 1.9]{Nguyen16}.

\section{Envelopes and comparison}
\label{sec:history-envelopes}

The semipositive big case of \cref{thm:envelope-contact} is \cite[Theorem 2.2]{ALS26}. The proof here gives the bounded-to-unbounded passage without bigness or a uniform mass assumption.

The contact conclusion in \cref{thm:envelope-contact-big} is proved under bounded mass in \cite[Theorem 2.9 and Proposition 2.11]{BGL25}. The proof in \cref{thm:big-general-envelope-contact} uses local replacement and a strict singular barrier, without that global mass assumption.

The fixed-point estimate in \cref{prop:fixed-point-finite} follows the argument of \cite[Proposition 2.12]{ALS26}. Restricting to finite sublevels avoids subtraction of possibly infinite measures.

The envelope proof of bounded domination in \cref{thm:semipositive-bounded-domination} is due to Guedj--Lu \cite[Proposition 2.8]{GL22}.

The difference-envelope construction used in \cref{def:relative-envelope-class} is recorded for semipositive big backgrounds in \cite[Propositions 2.4 and Lemma 2.5]{ALS26} and for arbitrary big backgrounds in \cite[Section 2.1]{XL26}. Nonemptiness alone gives a genuine envelope through the finite-sublevel bound.

The relative envelope class of \cref{def:relative-envelope-class} follows \cite[Definition 3.2]{ALS26} in the semipositive big setting.

The convex-combination argument in \cref{lem:relative-stretch-closure} is that of \cite[Proposition 3.5]{ALS26}.

The model condition in \eqref{eq:als-model-condition} and its individual positivity consequence appear in \cite[Proposition 2.10 and the discussion after Definition 2.13]{ALS26}. That paper does not assume bounded mass for its same-type equations with $L^p$ data. Individual positivity must be distinguished from a positive infimum over the whole ordinary type.

The lower-volume mechanism in \cref{thm:relative-uniform-mass-lower} is that of \cite[Proposition 3.8]{ALS26}.

The bounded-obstacle support argument in \cref{prop:lsc-obstacle-support} follows the approximation structure of \cite[Theorem 2.3]{GL22}.

The increasing bounded convergence used in \cref{prop:lsc-obstacle-support} is \cite[Theorem~7.4]{BT82}. Smooth local corrections and expansion against the remaining smooth background give its nonclosed-background application.

The modified comparison principle in \cref{thm:modified-comparison} is due to Ko\l odziej--Nguyen \cite[Theorem 0.2]{KN15}. The proof follows the finite induction of Guedj--Lu \cite[Theorem 1.5]{GL22}, with enlarged dimensional constants.

\section{Smooth and continuous equations}
\label{sec:history-smooth-equations}

For $n\geq2$, \cref{thm:tw-smooth} is \cite[Corollary 1]{TW10}. Their equation uses $\mathrm i\partial\bar\partial\varphi$; putting $u=2\pi\varphi$ gives the convention of this book. A general smooth positive $\mathrm{d}V$ is absorbed by replacing their prescribed function by $F+\log(\mathrm{d}V/\omega^n)$.

The weighted integration proof of the potential estimate in \cref{chap:smooth-equations} follows Cherrier and Tosatti--Weinkove \cite{Cherrier87,TW10}. A finite recurrence absorbs torsion before Sobolev iteration.

The Gauduchon concavity argument for elliptic normalization constants in \cref{thm:b-derivatives} is also used in \cite[Theorem 2.13]{Collins21}.

The prescribed Gauduchon-volume equation in \cref{thm:stw-volume} is \cite[Theorem 1.4]{STW17}. The potential in this book is $2\pi$ times the reference potential.

The weak prescribed-singularity results of \cite{ALS26,XL26} use different mass and reference assumptions, separated in \cref{chap:prescribed-equations,chap:relative-finite-measures,chap:bounded-mass-equations}. Uniqueness of the scalar and of the potential are distinct conclusions.

The continuous normalized equation in \cref{thm:kn-continuous} is \cite[Theorem 0.1]{KN15}. It requires neither bounded mass nor positive lower volume.

The continuous exponential equation \eqref{eq:weak-exp} is \cite[Theorem 0.1]{Nguyen16}. Both this paper and the preceding normalized-equation reference use $\ddc=\mathrm i\partial\bar\partial/\pi$, twice the operator here. Replace the reference potential $\varphi$ by $u=2\varphi$ and, for the exponential equation, use reference parameter $2\lambda$.

The bounded exponential comparison of \cref{lem:exp-order} recovers \cite[Lemma 2.3]{Nguyen16}. Neither density needs a positive lower bound.

The vanishing-exponent characterization of the normalization constant through \eqref{eq:exp-barriers} is due to Nguyen \cite[Corollary 0.2]{Nguyen16}. It establishes uniqueness of the scalar without proving uniqueness of the zero-exponent potential.

The transition from continuous to prescribed-singularity equations uses the relative-envelope frameworks of \cite{ALS26,XL26}. Their separate assumptions are recorded in \cref{sec:history-prescribed}.

The torsion calculation for the second-order estimate in \cref{chap:smooth-equations} goes back to Cherrier. Its completed-square form and the maximum argument follow \cite[Sections 3 and 9]{TW15}.

The Evans--Krylov input for \cref{lem:stw-higher-regularity} is \cite[Definition~3.2 and Theorem~3.4]{TWWY15}, for a concave uniformly elliptic operator with H\"older dependence on the base point.

The potential bound in \cref{sec:stw-czero} follows \cite[Propositions~10--11 and Remark~12]{Sze18}, with the determinant estimate specialized to the Gauduchon ansatz.

The rescaling step \eqref{eq:stw-rescaled-matrix} follows \cite[Section~6]{Sze18}. The first-order term is retained before rescaling and vanishes at rate $A_j^{-1}$.

The Euclidean Liouville proof in \cref{sec:stw-liouville} follows \cite[Theorem~20]{Sze18}, specialized to the cone arising from the Gauduchon equation.

The second-order proof in \cref{sec:stw-second-order-proof} specializes \cite[Section~3]{STW17} to the logarithmic determinant and retains the first-order terms.

The torsion identities \eqref{eq:stw-diagonal-torsion} are those used in \cite[Section~4]{STW17}. They explain the structure required of the first-order perturbation.

The local capacity deductions used for the weak equations recover the estimates of \cite[Section 5]{KN15} and \cite[Section 6]{Dinew16}. Their inputs are local pluripotential estimates, rather than an assumed global Hermitian capacity bound.

The uniform compactness statement \cref{thm:lp-uniform-compactness} is \cite[Corollary 5.10]{KN15}. The common density bound is essential; normalized $L^1$ compactness alone is insufficient.

The weighted mean \eqref{eq:gauduchon-exponential-mean} is Nguyen's mechanism \cite[Proposition 2.5]{Nguyen16}. The proof here obtains its lower bound directly from the sublevel estimate.

\section{Uniqueness and degeneracy}
\label{sec:history-uniqueness}

The restricted-density mass-concentration argument around \eqref{uni-eq:mixed-bounded} follows \cite[proof of Theorem~3.7 and Remark~3.8]{LPT20}. The proof makes its Hermitian comparison inputs explicit.

The exterior-ball proof leading to \eqref{uni-eq:normalzero} follows the maximum-propagation method of \cite[Proposition~2.4 and Theorem~2.5 in the cited preprint]{BG18}.

\section{Equations with prescribed singularities}
\label{sec:history-prescribed}

The general-measure conclusions of \cref{thm:li-general-measure} are \cite[Theorems 6.1--6.2 and Corollary 6.1]{XL26}. That paper assumes both metric volume properties. Under bounded mass, a model reference with positive lower volume implies metric positive lower volume, so the formulation here reorganizes that regime rather than enlarging it.

The domination results \cref{thm:relative-finite-domination,thm:big-relative-domination} use the finite-measure argument of \cite[Theorem 3.14]{ALS26} and the big-background comparison of \cite[Proposition 3.2]{XL26}. The proof of \cref{thm:relative-extended-domination} removes finiteness of the compared measure; the big-background result uses the positive lower volume of the reference.

The same ordinary-type equations of \cref{thm:als-prescribed-lp} are \cite[Theorems 5.4 and 5.7]{ALS26}. Replacing the reference volume $\omega^n$ by $\mathrm{d}V$ multiplies the density by a smooth positive factor and preserves its $L^p$ class. A positive constant change of the Hessian normalization is absorbed by rescaling both the potential and the exponential parameter.

The relative potential estimate in \cref{thm:relative-linfty} reconstructs \cite[Lemma 5.5]{ALS26}.

The rooftop supersolution of \cref{lem:relative-lp-supersolution} is the construction in \cite[Lemma 5.6]{ALS26}.

The dominated-continuity statement in \cref{thm:relative-dominated-continuity} is the fixed-measure case of \cite[Lemma 4.2 and Theorem 4.1, condition (2)]{ALS26}. The common dominating measure and positive lower mass prevent collapse of the lower envelopes.

The semipositive general-measure equations in \cref{thm:als-general-measure} are \cite[Theorem 5.8 and Corollary 5.9]{ALS26}. Their solutions have finite mass; the result does not assert bounded difference from the reference for arbitrary measure data.

The finite-measure domination theorem in \cref{chap:relative-finite-measures} reconstructs the finite-$\mu_u$ case of \cite[Theorem 3.14]{ALS26}. A fixed finite measure prevents collapse; other members of the envelope class need not have finite mass.

The semipositive general-measure regime summarized in \cref{chap:research-program} follows \cite{ALS26}, with the smaller-background model condition and positive lower mass.

The weighted addition theorem \cref{thm:weighted-measure-addition} reconstructs \cite[Theorem 5.2]{ALS26}, with a localized passage to the equation.

The bounded subsolution and existence argument in \cref{lem:big-lp-subsolution,thm:big-bounded-exponential} follows \cite[Lemma 3.3 and Theorem 3.4]{GL23E}.

The strict-barrier comparison and minimal-singularity equations developed from \cref{sec:big-minimal-comparison} reconstruct \cite[Section 4]{BGL25}. That section explicitly drops the bounded mass assumption used earlier in the paper.

The minimal ordinary-type equations culminating in \cref{thm:big-minimal-normalized} reconstruct \cite[Theorems 4.6--4.7]{BGL25}. They give finite mass of the solution without a mass bound on the entire ordinary type.

The minimal-reference general-measure equations in \cref{sec:big-minimal-general-measures} reconstruct \cite[Theorems 5.2--5.3]{XL26}. The relative volume assumptions and their relation to the metric assumptions are stated at that construction.

The higher-moment estimate in \cref{sec:big-minimal-moments} reconstructs the concave-weight argument of \cite[Proposition 5.1 and Lemma 5.1]{XL26}, with the integrability condition on the fixed singular measure stated explicitly. The measure-smoothing argument in \cref{lem:measure-smoothing-integrals} is the general psh input in the proofs of \cite[Theorems 5.2--5.3]{XL26}; its bounded-carrier integrability estimate is recorded in \cref{lem:bounded-carrier-integrability}.

The $L^p$ ordinary-type conclusion of \cref{cor:big-model-lp-type} is \cite[Corollary 6.1]{XL26}. The relative estimate \cref{prop:big-model-moment-bound} starts in the envelope class and uses a uniform moment of order $m>2n+1$. The concave-weight method is related to \cite[Theorem 5.1]{XL26}, whose stated hypothesis already assumes ordinary reference type and whose moment exponent is $m>n$. That statement is not used to infer the ordinary type of an arbitrary envelope-class solution.

The model-reference constructions in \cref{sec:big-model-general-measures} follow \cite[Theorems 6.1--6.2]{XL26}.

The relative first-moment mechanism behind \eqref{eq:big-energy-max-bound} is \cite[Proposition 4.2]{XL26}. The proof here writes the optimization over equivalent reference potentials explicitly.

The big relative foundations in \cref{sec:big-relative-foundations} reconstruct the envelope and domination results of \cite[Sections 3--4]{XL26}.

The fixed-measure continuity argument in \cref{sec:big-relative-foundations} proves \cite[Theorem 4.2]{XL26}. Finite rooftops use only the envelope class; the mass interval prevents collapse and controls tails.

The positive-volume stretching lemma \cref{lem:positive-volume-stretch} follows the convex-envelope argument of \cite[Proposition 3.1]{XL26}, with the mixed-mass estimate written explicitly.

The contact argument for relative projection in \eqref{eq:big-projection-contact-tail} follows \cite[Proposition 3.3]{XL26}; the proof first constructs the requisite rooftop.

The unrestricted addition argument using \eqref{eq:unrestricted-addition-cutoff} reconstructs \cite[Theorem 6.3]{XL26}. An additional rooftop supplies a common lower bound for the cutoff passage.

The four-stage weighted-addition construction in \cref{sec:big-weighted-addition} reconstructs \cite[Proposition 6.1]{XL26}.

The local decomposition used in \cref{lem:local-bounded-measure-carrier} is \cite[Theorem 6.3]{Ceg98}, on a bounded hyperconvex domain in dimension at least two. For a non-pluripolar measure, the pluripolar summand vanishes. The global constructions in \cref{lem:bounded-measure-carrier,prop:relative-measure-carrier} patch these local representations.

The bounded subsolution theorem \cref{thm:bounded-subsolution} reconstructs \cite[Theorem 5.1]{ALS26}.

The extraction argument in \cref{thm:relative-fixed-measure-continuity} reconstructs the fixed-dominating-measure case of \cite[Lemma 4.2 and Theorem 4.1]{ALS26}. Positive lower mass prevents collapse; bounded upper mass is not used.

The common-envelope argument preceding \cref{cor:exponential-relative-class-closure} supplies the noncollapse needed in the first alternative of \cite[Lemma 4.2]{ALS26}. For approximants of the reference's ordinary type, that noncollapse also holds without a positive lower-volume hypothesis.

The minimal-reference construction in \cref{sec:general-measure-minimal} is the corresponding case of \cite[Theorem 5.8 and Corollary 5.9]{ALS26}.

The prescribed-reference proof in \cref{sec:general-measure-relative} follows the density approximations in \cite[Theorem 5.8 and Corollary 5.9]{ALS26}.

The torsion calculation in \cref{lem:moving-bounded-integrals} reconstructs \cite[Lemma 2.3]{KN22}.

The common-barrier continuity of \cref{cor:common-barrier-continuity} reconstructs the common ordinary-type case of \cite[Theorem 4.1]{ALS26}. The reference need not be an envelope fixed point; the smaller-background condition suffices.

The singular strict-domination theorem \cref{thm:singular-strict-domination} reconstructs \cite[Theorem 3.11]{ALS26}.

\section{Normalization and adjoint measures}
\label{sec:history-normalization}

Gauduchon measures in the differential study of elliptic normalization constants also appear in \cite{Collins21}. The singular proof of \cref{thm:relative-normalization-concavity} uses mixed products and domination.

The selection theorem \cref{thm:vanishing-exponential-selection} reconstructs the supplied manuscript \cite[Theorem A, Propositions 3.1--4.1 and Lemma 5.1]{SuppliedSelection26}. The critical scalar defines a lattice norm on density roots; its norming face gives the probability family that normalizes the greatest critical subsolution. The subsolution envelopes then identify every cluster limit. The comparison and existence inputs are proved in \cref{sec:big-minimal-comparison,sec:big-minimal-existence}; \cref{thm:big-minimal-holder} gives the required common local modulus from the uniform reference interval and the bound on the actual equation densities. The scope is a fixed nonzero $L^p$ density, $p>1$, and ordinary minimal singularities in a smooth big background. This selects one critical solution without asserting uniqueness of all normalized critical solutions.

\section{Singularity classes and traces}
\label{sec:history-singularities}

The enlarged quotient of smooth real $(1,1)$-forms by smooth $\ddc$-terms in \cref{chap:nonclosed-bclasses} is used in \cite[Section 1]{BGL25}. Prescribing its exterior derivative retains the torsion on each modification.

The ruled-surface construction in \cref{sec:envelope-boundary-example} reconstructs \cite[Example 1.7]{DPS94} and applies it to recovery from envelope systems.

The approximation comparison used for \cref{thm:trace-data-additivity} is the resolution argument underlying Cao's addition theorem \cite[proof of Theorem~1.4]{Cao15}. It needs no global closed background.

\section{Sections and algebraic dimension}
\label{sec:history-sections}

The projective-modification form of Moishezon's theorem used in \cref{chap:flag-bodies} is recorded in \cite[Introduction, p.~1332]{Moi67}. It is retained as a classical geometric input.

The bounded-mass consequence of a K\"ahler modification in \cref{lem:kahler-modification-bounded-mass} is also the upper-volume pullback argument of \cite[Lemma 3.6]{GL22}.

For section rings constructed through algebraic reduction, see \cite{Wu24}. The flag construction in \cref{chap:flag-bodies} uses a specified compact flag on an auxiliary modification; it makes no assertion about centers of arbitrary meromorphic-field valuations.

The calculation and sharp example in \cref{sec:horizontal-center-constraints} reconstruct the supplied manuscript \cite[Theorem~1.1, Lemma~2.2 and Proposition~4.1]{SuppliedHorizontal26}. The chapter combines them with its kernel comparison.

\section{Gauduchon variational volume}
\label{sec:history-gaud-volume}

The concavity method underlying \cref{chap:gauduchon-variational-volume} belongs to the general construction of \cite[Theorem 2.13]{Collins21}. The chapter develops its Gauduchon formula and the role of the chosen probability measure.

Completeness and finite-distance positive-volume boundaries for the closed K\"ahler Hessian metric are studied in \cite[Proposition 4.4 and Lemma 4.5]{Magnusson12}. The homogeneous argument in \cref{sec:gauduchon-metric-geometry} uses the present normalization on a finite-dimensional cone.

On the closed K\"ahler cone, the metric in \cref{sec:gaud-volume-hessian} is $1/n$ times the Hessian metric of $-\log\int_X\theta^n$ from \cite[Section 1]{Magnusson12}.

The existence of a scalar conformal factor for a weak Monge--Amp\`ere power is raised in \cite[Chapter 3, Section 5]{LuHDR22}. The finite-current representation of \cref{thm:gaud-current-all-adjoints} supplies a supporting probability without defining that weak power or its scalar factor.

\section{Closed classes and singular loci}
\label{sec:history-closed-loci}

The notation for almost positive representatives in \cref{sec:bc-hermitian-bounds} follows Boucksom \cite[Section 2.1.2]{BouThesis02}. It expresses the lower bound without treating a nonclosed Hermitian form as a Bott--Chern class.

The minimal multiplicities in \cref{def:bc-minimal-multiplicity} are those of \cite[Definition 2.1.5]{BouThesis02}.

The smooth and analytic regularizations in \cref{subsec:psh-closed-regularization}, used for \cref{prop:bc-non-nef}, have the lower bounds in \cite[Theorem 2.1]{Bou04Zariski}; see also \cite[Theorems 1.2.9--1.2.10]{BouThesis02}. Smooth regularization loses a multiple of a continuous approximation of the Lelong numbers; analytic regularization has vanishing uniform curvature error and does not increase Lelong numbers.

The nef ruled-surface example in \cite[Section 1.2.3]{BouThesis02} has a positive current with minimal singularities and nonzero Lelong numbers. It explains why \cref{def:bc-minimal-multiplicity} uses almost positive approximating families.

The non-K\"ahler locus in \cref{def:bc-non-kahler-locus} is defined in \cite[Definition 2.1.19]{BouThesis02}.

The finite-intersection and regularized-maximum proof of \cref{thm:bc-non-kahler-locus} is the argument of \cite[Theorem 2.1.20]{BouThesis02}.

The equivalence between a K\"ahler current and a K\"ahler modification, used in \cref{chap:bott-chern-loci}, is the classical criterion recalled in \cite[Theorem 1.2.11]{BouThesis02}. It is retained as an external geometric input.

The equality of the non-K\"ahler and null loci for a nef and big class is the transcendental Nakamaye theorem of Collins--Tosatti \cite{CT15Loci}, also valid on compact manifolds in Fujiki's class $\mathcal C$. The discussion in \cref{chap:bott-chern-loci} recalls this external theorem.

The numerical restricted volume in \eqref{eq:bc-numerical-restricted-volume} follows Collins--Tosatti \cite[Section 2]{CT22Restricted}. The smooth remainder in the analytic singularities is part of that definition.

\section{Closed classes and metric completeness}
\label{sec:history-completeness}

The normal-length argument in \cref{sec:divisorial-normal-length} is adapted from \cite[Theorem~1.1]{SuppliedCompleteness26}. A Schur complement and an exhaustion function give the metric estimate before the exponential equation supplies its determinant bound.

The torus profiles in \eqref{eq:torus-length-profiles} are those of \cite[Lemma~3.1 and Propositions~3.2--3.3]{SuppliedCompleteness26}. Their global mass and metric properties are checked directly in the chapter.

The examples in \cref{sec:divisorial-normal-length} do not address the flow compatibility conditions of \cite[Section~3.1]{DNGL26}. Their densities have logarithmically weakened divisorial poles, rather than a nonintegrable pole with a smooth positive coefficient across the divisor.

The divisorial exponential completeness question tested by \cref{thm:incomplete-divisorial-metric} is raised in \cite[Question~3.5]{DNGL26}. The example concerns arbitrary big classes; flow-compatible data require further conditions.

The finite-energy class used in \cref{sec:incomplete-divisorial-metric} is the classical class in \cite[Section~2]{BEGZ10}, relative to a minimal-singularity reference in a big closed class.

The finite mixed masses and relative energy used in \cref{thm:incomplete-divisorial-metric} are those of \cite[Proposition~1.6 and Section~2.2]{BEGZ10}.

Uniqueness in \cref{thm:incomplete-divisorial-metric} uses the full-mass comparison principle \cite[Corollary~2.3]{BEGZ10}. Although the unweighted reference measure is infinite, the two compared Monge--Amp\`ere measures have finite mass.

The compatibility conditions in \cite[Section~3.1]{DNGL26} require the limiting class, divisor and singularity coefficients to arise from specified flow data. The incomplete big-class example does not settle that problem.

\section{Locally conformally K\"ahler geometry}
\label{sec:history-lck}

The balancedness conjecture stated in \cref{chap:lck} is due to Ornea--Verbitsky \cite{OV25Balanced}.

The proof of \cref{thm:lck-chern-ricci-rigidity} specializes the argument of \cite[Section 4]{BO25LCK} to the Chern--Ricci-flat case.

The prescribed-volume theorem in \cref{chap:lck} is \cite[Theorem 4.1]{OV24Vaisman}, including invariance of the volume form under both the Lee and anti-Lee vector fields.

The conjecture following \cref{thm:lck-chern-ricci-rigidity} comes from \cite{BO25LCK}. Vanishing first Bott--Chern class supplies a Chern--Ricci-flat conformal representative, which need not be the Gauduchon representative required by the rigidity argument.

The dual description and bimeromorphic invariance in \cref{sec:lee-gauduchon-cone} are due to Ornea--Verbitsky \cite[Theorems 3.11 and 4.5]{OV25LeeGauduchon}.

The distinction between vanishing Lee--Gauduchon class and the strongly Gauduchon condition in \cref{sec:lee-gauduchon-strong} follows \cite[Remarks 5.6--5.8]{OV25LeeGauduchon}.

The strongly Gauduchon existence criterion in \cref{sec:lee-gauduchon-strong} is due to Popovici \cite[Proposition~1.2]{Pop13SG}.

\section{Volumes under mass bounds}
\label{sec:history-volumes}

The singularity volumes used in \cref{sec:bounded-mass-volume-comparison} are those of \cite[Definition 3.1]{BGL25}, expressed in potentials.

The closed-background identity for the two singularity volumes is recalled in \cite[Example 3.2]{BGL25}. It uses invariance of non-pluripolar mass under bounded changes of singularities.

The proof of \cref{thm:volume-monotonicity} reconstructs \cite[Proposition 3.7]{BGL25} in the big regime.

The lower-volume addition argument in \cref{prop:volume-addition} is the envelope proof of \cite[Proposition 2.5]{XL26}.

The relative full-mass class in \cref{def:relative-volume-class} is \cite[Definition 3.2]{XL26}.

The collapse argument in \cref{thm:relative-equivalence} is the full-volume case of \cite[Lemma 3.1]{XL26}. It uses bounded mass and the positive lower volume of the reference.

The implication from positive lower volume to the smaller-background condition in \cref{prop:positive-volume-smaller-background} is the forward direction of \cite[Proposition 2.11]{ALS26}. The proof bounds its mixed-product error explicitly using bounded mass. This direction requires neither the ordinary envelope identity nor $v_-(\theta)>0$; the latter assumption is used for the converse. Applying the implication to an ordinary envelope fixed point compares \eqref{eq:li-model-condition} with \eqref{eq:als-model-condition}.

The relative-envelope results of \cite{ALS26,XL26} address parts of the modelization problem in \cref{qu:hermitian-i-characterization}. The remaining questions should not be read as claims that these existing results are absent.

\section{Regularity and singular spaces}
\label{sec:history-regularity}

For Hermitian backgrounds, Ahn \cite[Theorems 1.1--1.2]{Ahn26} proves locality of H\"older regularity and regularity of weighted extremals under a uniform density-in-capacity condition on the compact set. These are the specific regularity regimes referred to in \cref{chap:research-program}.

Dang \cite[Theorem B]{Dang26} proves continuity results for bounded solutions on resolutions under additional geometric and integrability hypotheses, notably the log terminal setting. This result does not cover arbitrary unibranch singularities in the problem posed in \cref{chap:research-program}.

Pan \cite[Theorems A and A']{Pan25} constructs a Gauduchon metric on the regular locus of a compact reduced irreducible Hermitian space of dimension at least two. It is conformal and uniformly comparable to the given metric, and unique up to scale within that class. Control of conformal factors on a resolution is central to the construction. This external result motivates the singular trace questions in \cref{chap:research-program}.

Existence and uniqueness for \eqref{eq:big-class-smooth-volume} are \cite[Theorem~A, equivalently Theorem~3.1]{BEGZ10}; the nef case of its ample-locus smoothness conclusion is \cite[Theorem~C]{BEGZ10}. The proof here treats smoothness without nefness.

The minimal bounds \eqref{eq:ample-solution-minimal-bounds} are the classical estimate \cite[Theorem~B, quantitatively Theorem~4.1]{BEGZ10}.

The local H\"older input used for the ample-locus argument is reconstructed in \cref{thm:big-minimal-holder} from \cite[Theorem~3.1]{Dang26}.

The regularity theorem \cref{thm:big-minimal-holder} reconstructs the local H\"older conclusion of \cite[Theorem~A and Theorem~3.1]{Dang26}. The proof keeps the unbounded global potential and nonclosed lower curvature bound explicit.

The formulation of \cite[Theorem~3.1]{Dang26} with a positive scalar and an equation only on the regular locus follows from \cref{thm:big-minimal-holder}: the analytic complement carries neither smooth volume nor non-pluripolar mass.

The strict-barrier stability estimate in \cref{sec:big-reference-stability} reconstructs \cite[Theorem~3.3 and Remark~3.4]{Dang26} in the minimal case, with the normalization of the comparison potential explicit.

The auxiliary-subsolution method of \cref{sec:bounded-stability} is also used for regularity in \cite[Theorem 3.3 and Remark 3.4]{Dang26}. The proof here first treats a semipositive big background.

The local boundary criterion and examples of \cref{sec:local-density-boundary} reconstruct the supplied manuscript \cite[Theorems~1.1--1.2 and Sections~3--5]{SuppliedLocalDensity26}.

The smooth Dirichlet input for \eqref{eq:local-density-determinant} is \cite[Theorem~1.1, p.~1186]{GL10}, used with zero background, flat metric $(\mathrm{i}/2)\partial\bar\partial|z|^2$ and determinant density $G$. A large quadratic correction supplies the strict subsolution.

The local identification in \cref{sec:local-density-boundary} uses bounded Dirichlet uniqueness \cite[Corollary~4.4, p.~17]{BT82}.

The smooth positive-density Dirichlet input for \cref{sec:local-homogeneous-dirichlet} is Guan--Li \cite[Theorem 1.1]{GL10}. The section proves the subsequent passage to continuous boundary values and zero density.

The variable-coefficient Evans--Krylov input for \eqref{eq:interior-local-ma} is \cite[Definition~3.2, Theorem~3.4, and the paragraph following it]{TWWY15}. Once uniform ellipticity holds, difference quotients and Schauder estimates give smoothness.

The smooth local approximants in \cref{sec:interior-approximants} use \cite[Theorem~1.1]{GL10}, with a strict quadratic subsolution. In dimension one the input is the Poisson Dirichlet problem.

The stability theorem \cite[Theorem 4.1]{KN21} assumes $f\geq c_0>0$ in addition to capacity hypotheses. The nonuniqueness densities in \cref{chap:degenerate-uniqueness} have no such lower bound, even when they are positive almost everywhere.

\section{Geometric problems}
\label{sec:history-geometry}

The K\"ahler-current characterization of Fujiki's class used in \cref{sec:non-kahler-geometric-problems} is \cite[Theorem 0.6 in the arXiv version]{DP04}. The transcendental Grauert--Riemenschneider conjecture is \cite[Conjecture 0.7 in the arXiv version]{DP04}. These locators refer specifically to the arXiv version.

Under nefness and positive top self-intersection, the equivalence between bigness and bounded mass cited in \cref{sec:non-kahler-geometric-problems} is \cite[Theorem 4.9]{GL22}. The characterization of Fujiki's class by bounded mass and positive upper class volume is \cite[Theorem C]{BGL25}; that upper volume differs from the fixed-probability Gauduchon functional.

The non-Hermitian locus in \cref{sec:non-kahler-geometric-problems} uses \cite[Definition 2.12 and Theorem 2.14]{DangNull26}. Its almost analytic singularities have the exact meaning of \cite[Definition 2.8]{DangNull26}, including a smooth remainder on a suitable log resolution.

Dang's qualitative null-locus equality for nef and big, possibly nonclosed backgrounds is \cite[Theorem 1.1]{DangNull26}. Theorem A of that paper gives a numerical criterion under bounded mass on the ambient manifold and positive lower volume on its positive-dimensional submanifolds. These are distinct scopes.

Theorem C of \cite{DangNull26} identifies finite-time Chern--Ricci singularities with the null locus under positive limiting lower volume, for arbitrary initial metrics in dimension two and Guan--Li initial metrics in higher dimensions. The latter satisfy $\ddc\omega_0=0$ and $\mathrm{d}\omega_0\wedge\dc\omega_0=0$.

Once the limiting background is big, \cite[Theorem 4.3]{DangNull26} gives current convergence and smooth local convergence off its non-Hermitian locus. This is the parabolic input distinguished from the numerical bigness question in \cref{qu:chern-ricci-gaud-singularities}.

The geometric smoothing conditions in \cite[Section~3.1]{DNGL26} remain separate from the unrestricted big-class counterexample used in \cref{sec:non-kahler-geometric-problems}.

The balanced--pluriclosed conjecture referred to in \cref{sec:non-kahler-geometric-problems} is stated in \cite[Conjecture 1.1]{FV26}. The balanced and pluriclosed metrics are allowed to be different.

The assertion that a pluriclosed manifold in Fujiki's class is K\"ahler is \cite[Theorem 0.2 in the author manuscript]{Chiose14}. The locator refers to the author manuscript used in \cref{sec:non-kahler-geometric-problems}.

The proposed long-time existence and convergence of pluriclosed flow in the balanced setting with $c_1(X)=0$ is \cite[Conjecture 1.2]{FV26}. The discussion in \cref{sec:non-kahler-geometric-problems} distinguishes a conserved balanced trace from a pointwise metric bound.

The closed K\"ahler transcendental Morse inequality, volume differentiability and cone-duality conclusions mentioned in \cref{sec:non-kahler-geometric-problems} are \cite[Theorem 1.1 and Corollary 1.2]{Tosatti26}, retained as literature results without reconstructing their proof. The qualitative criterion under bounded mass is \cite[Theorem 4.12]{GL22}.

The maximal-time theorem used in \cref{sec:smooth-cone-current-duality} is \cite[Theorem 1.2 and (1.4), arXiv version]{TW15}. Smooth potentials are multiplied by $2\pi$ in the convention of this book; time is unchanged because the Ricci form is the one used in that source.
